\documentclass[11pt,a4paper]{article}

\usepackage[utf8]{inputenc}
\usepackage[T1]{fontenc}
\usepackage{lmodern}
\usepackage{amsmath,amssymb,amsthm}
\usepackage{booktabs}
\usepackage{enumitem}
\usepackage{hyperref}
\usepackage[margin=2.5cm]{geometry}
\usepackage{xcolor}

\definecolor{codeblue}{RGB}{30,90,156}

\newtheorem{theorem}{Theorem}[section]
\newtheorem{proposition}[theorem]{Proposition}
\newtheorem{corollary}[theorem]{Corollary}
\newtheorem{lemma}[theorem]{Lemma}

\theoremstyle{definition}
\newtheorem{definition}[theorem]{Definition}
\newtheorem{claim}[theorem]{Claim}

\theoremstyle{remark}
\newtheorem*{remark}{Remark}
\newtheorem*{boundary}{Boundary}

\hypersetup{
  colorlinks=true,
  linkcolor=codeblue,
  citecolor=codeblue,
  urlcolor=codeblue,
}

\newcommand{\Ndelta}{\mathbb{N}_{\delta}}
\newcommand{\Zdelta}{\mathbb{Z}_{\delta}}
\newcommand{\Qdelta}{\mathbb{Q}_{\delta}}
\newcommand{\Qpos}{\mathbb{Q}_{>0}}
\newcommand{\Rpos}{\mathbb{R}_{>0}}

\title{%
  A Single Primitive for Mathematics and Physics:\\
  What Distinction Forces, What It Posits, and Where the Boundary Lies}
\author{%
  Jonathan Washburn\\
  Recognition Physics Institute, Austin, TX, USA\\
  \texttt{jon@recognitionphysics.org}}
\date{May 2026}

\begin{document}

\maketitle

\begin{abstract}\sloppy
This paper is the synthesis of a program that takes a single primitive, the act
of distinction, and asks how far it determines the structures of mathematics and
physics. We give the honest map. From one act of distinction, with no axiom of
choice, no principle of excluded middle, and no axiom of infinity, the natural
numbers, the integers, and a field of ratios are constructed, each with its
arithmetic laws proved and the strength of each proof recorded. The canonical
reciprocal cost that the physical program rests on is \emph{not} determined by
that field of ratios alone: the space of admissible costs is classified and the
obstruction is exhibited exactly, as a free choice of orientation on each prime.
The continuous completion of the ratio field, the one structure on which the cost
becomes forced, is itself \emph{not} a consequence of distinction; it is a
separate, named commitment, and the gap is a gap in cardinality, the countable
reach of distinction against the uncountable continuum. On that completion, and
only there, four algebraic laws force the cost into a single one-parameter family
that is exactly the gauge orbit of the canonical cost, and one calibration datum
selects the canonical representative uniquely within the orbit. The value of that
datum is, at present, a posit rather than a consequence. The resulting picture is
stratified, not seamless: mathematics is what distinction permits, physics is
what the completion of distinction costs, and the two meet at one named boundary,
the passage from the countable ratio field to its uncountable completion. We state
the claim at exactly this strength, mark each link as forced, posited, or open,
and name the single research result, deriving the calibration unit from the cost
of one distinction act, that would convert the cost half of the program from a
gauge classification into an unconditional forcing.
\end{abstract}

\tableofcontents

\section{Introduction: one primitive, an honest map}

A foundational program that proposes to derive both mathematics and physics from a
single primitive owes two things to its reader. It owes a precise statement of the
primitive, and it owes an honest accounting of how far the primitive actually
reaches: what it forces, what it merely permits, what it must posit, and where it
runs out. The temptation, in a program of this ambition, is to present an
unbroken chain from the primitive to the constants of physics and to call the
chain a derivation. We resist that presentation here, because it is not true, and
because the truth is both more defensible and more interesting.

The primitive is the act of distinction: the drawing of a difference, the marking
of one thing as not another. We write it $\delta$. The claim of the program is not
that $\delta$ forces everything downstream in one seamless sweep. The claim, stated
at the strength we can defend, is a stratification in three layers with one
boundary between them.

\begin{quote}
\emph{From a single act of distinction, with no axiom of choice, no excluded
middle, and no axiom of infinity, the natural numbers, the integers, and a
rational field are constructed, each with its arithmetic laws proved and its
proof strength audited. The canonical reciprocal cost $J$ is not determined by
this rational carrier alone; the space of admissible costs is classified, and the
obstruction is exhibited exactly. On the continuous completion, a named
commitment that is itself not forced by distinction, shown by a cardinality
argument, the four algebraic laws force exactly the one-parameter gauge orbit of
$J$, and a single calibration datum selects $J$ uniquely within it. Mathematics
is what distinction permits; physics is what the completion of distinction costs;
the two meet at one named boundary.}
\end{quote}

Everything in this paper is in service of that paragraph, and nothing in this
paper exceeds it. The three technical results it summarizes are developed in full
in three companion papers: the construction of the number tower from
distinction~\cite{DeltaArithmetic}, the proof that distinction does not force the
continuum~\cite{DeltaContinuum}, and the classification of the cost as a forced
form with a gauge unit~\cite{DeltaCost}. This synthesis states the terminal claim,
assembles the three results into a single map, makes the strength of each link
explicit, treats the question of inevitability honestly and conditionally, and
names the one open frontier whose resolution would strengthen the whole.

\subsection{Why stratified, not seamless}

The honest map has a shape, and the shape matters. The first layer, arithmetic, is
forced: distinction alone, with the weakest possible logical commitments,
produces the whole number tower. The second layer, the continuum, is \emph{not}
forced: distinction reaches only countably far, and the continuous line that
physics requires is a strictly larger object that must be adopted as a separate
commitment. The third layer, cost, is forced \emph{only on the completion} and
\emph{only up to a gauge}: on the countable ratio field the cost is wide open, on
the completion the algebraic laws cut it to a single orbit, and one unit of scale
selects the canonical member.

The boundary between layer one and layer three is therefore not a technicality. It
is the exact place where the program crosses from mathematics into physics, and it
coincides with the place where distinction stops forcing and starts positing. That
the two coincide is the central structural fact of the program, and we will show
why: the cost is underdetermined on the ratio field for the same reason the ratio
field is not complete, namely that a gapped domain insulates local constraints
from global ones. Remove the gaps, complete the field, and the same act that fixes
the domain also fixes the cost form. The price of removing the gaps is a posit.
The honest program names that price.

\subsection{What this paper does not claim}

Three claims are deliberately absent, and a reader familiar with the program's
earlier and more exuberant presentations should note their absence as deliberate.

First, we do not claim that distinction forces the canonical cost $J$. On the
ratio field it does not, and this is a theorem about the ratio field, not a gap
waiting to be filled. We claim that distinction forces the cost \emph{form}, the
gauge orbit, and that one datum fixes the gauge.

Second, we do not claim that distinction forces the continuum. It does not, and
this too is a theorem. The completion is a posit, and we say so in the abstract,
in the introduction, and again at the boundary where it is invoked.

Third, we do not claim that the calibration unit, the one datum that selects $J$
within its orbit, is itself derived from distinction. At present it is a posit, a
choice of scale. Whether it can be derived is the open frontier of
Section~\ref{sec:frontier}; until it is, claiming it would be claiming more than is
proved.

With these absences marked, the positive content is large, and it is exact. We
turn to it.

\section{The primitive}\label{sec:primitive}

\subsection{Distinction, and the honest question of how many primitives}

The act of distinction is the marking of a difference. To distinguish is to hold
two things and judge them not the same. From this single act the program builds
everything that follows, and the first duty of an honest account is to say
exactly what the act contains.

It contains two moves that we must not blur. There is the act of marking, the
production of a new token, and there is the judgment that compares two tokens and
returns same or different. The program's strongest slogan would be that these are
one primitive, that the comparison is internal to the act. We do not assert that
here, because the construction as it stands uses both: an operation that produces
a successor token, and a decidable judgment of sameness on tokens. Whether the
judgment reduces to the act, so that distinction is genuinely one primitive, or
whether recognition is irreducibly the pair, mark and compare, is an open
question that we state plainly in Section~\ref{sec:frontier}. The honest position
is that recognition is at most two primitives and possibly one, and that nothing
downstream depends on resolving this, because both the marking and the comparison
are available from the start.

What matters for the construction is that the judgment is \emph{decidable}: for
any two tokens, same or different is determined, with no appeal to the law of
excluded middle as a logical axiom, because the determination is computed, not
asserted. This is the first of the strength commitments we track.

\subsection{Strength tags}

The single discipline that holds the whole program honest is that every result
carries a tag recording the strength of commitment its proof requires. We use
three tags, in increasing order of commitment.

\begin{itemize}[leftmargin=2em]
  \item \textbf{Distinction-only.} The result follows from the act of
    distinction and finite repetition alone: marking, comparing, and iterating a
    finite number of times. No completed infinite totality, no choice principle,
    no nonconstructive case split. This is the weakest and most valuable tag.
  \item \textbf{Trace-closure.} The result requires a completed family of traces,
    an infinite totality treated as a finished object, such as the set of all
    finite orbits or a limit of a sequence of them. This is a genuine strengthening
    over distinction-only, and the completion of the ratio field will carry this
    tag.
  \item \textbf{Classical extension.} The result requires a classical principle,
    excluded middle on an undecidable predicate, or a choice principle, used
    essentially rather than as mere verifier convenience.
\end{itemize}

The tags are not decoration. They are the content of the honesty claim. When we
say arithmetic is forced and the continuum is not, what we mean precisely is that
the number tower carries the distinction-only tag throughout, and the completion
carries the trace-closure tag and cannot be brought down to distinction-only. The
boundary between mathematics and physics in this program is, exactly, the boundary
between the distinction-only tag and the trace-closure tag.

\section{What distinction permits: the number tower}\label{sec:arithmetic}

The first layer of the map is the arithmetic that distinction forces. We summarize
the construction and its results; the full development with proofs is the subject
of the companion arithmetic paper~\cite{DeltaArithmetic}.

\subsection{Natural numbers}

Iterating the act of distinction produces a sequence of tokens, each marked off
from its predecessor. The orbit of a token under repeated marking is a
distinction-native natural number; we write the type $\Ndelta$. A successor
operation and a zero are immediate, the decidable sameness judgment gives equality,
and induction is available because every token is reached from zero by finitely
many marks. On $\Ndelta$ one defines addition and multiplication by the usual
primitive recursions, and the semiring laws, commutativity, associativity,
distributivity, the identities, follow by induction. The order is decidable and
agrees with the arithmetic.

\begin{claim}[Arithmetic of $\Ndelta$, distinction-only]\label{cl:nat}
The distinction-native naturals $\Ndelta$ form a commutative semiring with
decidable equality and a decidable linear order, and the structure is isomorphic
to the standard naturals. Every law is proved under the distinction-only tag: no
completed infinity, no choice, no nonconstructive case split.
\end{claim}

The content of Claim~\ref{cl:nat} is not that $\Ndelta$ is a novel object; it is
isomorphic to the naturals, as it must be. The content is the \emph{tag}. The
standard set-theoretic construction of the naturals invokes an axiom of infinity to
collect them into a completed set. The distinction-native construction does not:
each natural is a finite object, the type is open-ended rather than a completed
totality, and no step appeals to a totality of all naturals. That is a strictly
weaker commitment, and it is the first thing the program can say that ordinary
foundations cannot.

\subsection{Integers and ratios}

Signed orbits, pairs of natural orbits read as a difference modulo the usual
cancellation, give the integers $\Zdelta$, with the full commutative ring
structure and a decidable order extending that of $\Ndelta$, all under the
distinction-only tag. Ratio orbits, pairs of a signed orbit and a nonzero natural
orbit read as a quotient modulo cross-multiplication, give a field $\Qdelta$ of
ratios, with the field operations, the decidable order, the absolute value, and
the reciprocal, again distinction-only. The field $\Qdelta$ is the
distinction-native rational field, and it is isomorphic to the rationals.

\begin{claim}[The rational field, distinction-only]\label{cl:rat}
The distinction-native ratios $\Qdelta$ form an ordered field with decidable
equality and order, isomorphic to the rational field, and every law is proved
under the distinction-only tag.
\end{claim}

On this field the permission structure of arithmetic, divisibility, primality, the
Euclidean algorithm, greatest common divisors, is available and decidable. We call
this the \emph{permission} reading of distinction: it tells us what comparisons of
orbits are allowed, which orbit divides which, which orbit is prime. Permission is
qualitative. It is the mathematics that distinction permits, and it is forced, all
of it, under the weakest tag.

\begin{boundary}
The number tower stops at $\Qdelta$. The next object the physical program needs,
the golden ratio, is a root of $x^2 = x + 1$, hence involves $\sqrt 5$, which is
not in $\Qdelta$. So even the first downstream physical quantity already lies
outside the distinction-native field. This is the first signal that physics does
not live on the ratio field, and it points directly at the boundary of
Section~\ref{sec:continuum}.
\end{boundary}

\section{Where necessity ends: the continuum is not forced}\label{sec:continuum}

The second layer of the map is a negative, and it is the hinge of the whole
program. We summarize it; the full argument is the companion continuum
paper~\cite{DeltaContinuum}.

The physical program needs the continuous real line. It needs it because the cost
becomes forced only there (Section~\ref{sec:cost}), because the golden ratio and
the other downstream constants live there and not on $\Qdelta$, and because the
limiting and second-derivative notions that physics uses have no traction on a
gapped field. The question is whether distinction forces this continuum, the way
it forces the number tower. It does not, and the reason is a counting argument
that no amount of cleverness in the construction can evade.

\begin{claim}[The continuum is not forced]\label{cl:cont}
The objects that distinction generates form a countable collection: the tokens,
the orbits, the signed orbits, the ratios, and any family of these produced by
finitely-branching generative rules iterated countably, are all countable. The
continuous real line is uncountable. Therefore distinction, whose entire output is
countable, cannot force the continuum. The gap between them is a gap in
cardinality: the countable reach of distinction against the uncountable line. The
completion that closes the gap is a separate commitment, carrying the
trace-closure tag, and it is provably stronger than anything distinction-only can
reach.
\end{claim}

The mathematics behind Claim~\ref{cl:cont} is elementary: it is Cantor's theorem
that the reals are uncountable, together with the observation that a generative
primitive producing finitely much at each of countably many stages reaches only a
countable total. The companion paper makes both halves precise and adds a second,
independent argument from model theory, that a countable first-order description of
distinction has a countable model, so no such description pins the continuum. The
contribution is not the diagonal argument. The contribution is the placement of
the boundary: showing that this specific generative primitive, the one the program
is built on, cannot reach the continuum, and therefore that the continuum is an
addition to the program rather than an output of it.

This is where the program crosses from mathematics into physics, and we mark the
crossing precisely. Up to and including $\Qdelta$, everything is
distinction-only, and everything is forced. The completion to the real line is
trace-closure, and it is posited. Physics begins at the posit.

\begin{boundary}[The completion is named, not derived]
We adopt the continuous completion of $\Qdelta$ as an explicit commitment. We do
not derive it from distinction, and we do not pretend to. Every result downstream
of this point that requires the completion inherits the trace-closure tag, and we
carry that tag forward honestly rather than laundering it into a distinction-only
claim.
\end{boundary}

\section{What the completion costs: the cost form forced, the unit a gauge}\label{sec:cost}

The third layer is the cost, and it is the load-bearing joint of the entire
program, because almost everything in the physical theory downstream flows from
the precise form of the cost. We summarize the result; the full treatment, negative
first, is the companion cost paper~\cite{DeltaCost}, which in turn imports the
continuous uniqueness theorem from the functional-equation
work~\cite{LogicFE,AxiomsCost}.

The cost is the canonical reciprocal cost
\[
  J(x) = \tfrac{1}{2}\bigl(x + x^{-1}\bigr) - 1,
\]
the price of comparing two magnitudes whose ratio is $x$. The structural
requirements a cost is asked to satisfy are reciprocal symmetry, $F(x) = F(x^{-1})$;
normalization, $F(1) = 0$; a composition law governing how costs combine; and,
when the domain is continuous, continuity and a unit calibration fixing the local
curvature at the unit. The question is whether these requirements force $J$. The
answer splits exactly along the boundary of Section~\ref{sec:continuum}.

\subsection{On the ratio field, the cost is not forced}

\begin{claim}[Discrete underdetermination]\label{cl:costneg}
On the distinction-native ratio field, the requirements of reciprocal symmetry,
normalization, and the composition law do not force $J$. The admissible costs are
exactly the character costs $F_\chi(u) = \tfrac{1}{2}(\chi(u) + \chi(u)^{-1}) - 1$,
where $\chi$ ranges over the homomorphisms of the multiplicative group of positive
ratios. Because that group is free abelian on the primes, a character is a free,
independent choice of value on each prime, and the admissible costs form an
infinite product with one factor per prime. The canonical cost $J$ is the single
character given by the inclusion; for every prime there is an admissible cost that
agrees with $J$ on all other primes and inverts the orientation of that one.
\end{claim}

Claim~\ref{cl:costneg} is the honest negative, and it is structural rather than
anecdotal: the obstruction is not a handful of counterexamples but the entire
character group. A uniqueness statement that recovers $J$ on the ratio field must
first fix the cost on every prime, which is to assume agreement with $J$ on a
generating set; that is a fact about generation, not a derivation of $J$ from
distinction. We never present generator-agreement as forcing.

\subsection{On the completion, the cost form is forced}

\begin{claim}[Continuous form forcing and the gauge orbit]\label{cl:costpos}
On the continuous completion, continuity restricts characters to power maps, and
the four laws force the cost into the one-parameter family
\[
  F_\lambda(x) = \tfrac{1}{2}\bigl(x^{\lambda} + x^{-\lambda}\bigr) - 1,
  \qquad \lambda > 0,
\]
which is exactly the solution set, not merely a subset of it. This family is the
orbit of $J$ under the group of multiplicative automorphisms of the positive
reals, and the action is free and transitive, so the family is a torsor: the
residual freedom is exactly one positive real, a gauge coordinate. A single
comparison value pins the gauge, and the unit calibration, log-curvature one at the
unit, selects the canonical member $F_1 = J$.
\end{claim}

The two claims together are the precise sense of ``the cost form is forced, the
unit is a gauge.'' Distinction and the algebraic laws force the form, the gauge
orbit. One datum fixes the gauge. The value of that datum is a
multiplicative-automorphism gauge that distinction does not, by itself, fix. The
entire arbitrary content of the cost layer is exactly one positive real number, and
everything else about the cost is forced.

\begin{boundary}[The calibration unit is a posit, for now]
The unit calibration is supplied from outside, as a choice of scale. We do not
claim it is derived from distinction. Whether it can be is the open frontier of
Section~\ref{sec:frontier}.
\end{boundary}

\subsection{Sharpening: the cost form needs only order, not the continuum}\label{sec:cost-monotone}

Claim~\ref{cl:costpos} was stated with continuity, the hypothesis under which a
multiplicative endomorphism of the positive reals is a power map. Continuity is
stronger than the argument requires. A monotone additive function is linear
(Darboux), so a \emph{monotone} multiplicative endomorphism is already a power map,
and the form forcing of Claim~\ref{cl:costpos} goes through with ``continuous''
replaced by ``order preserving.''

The consequence reshapes the boundary. Continuity is about limits, and limits need
the gaps of $\Qdelta$ filled, which is the move to the uncountable line.
Monotonicity is about order, and order is already present on the real-algebraic
closure of $\Qdelta$, the countable field in which every distinction-posable
polynomial comparison resolves. That field contains $\sqrt{2}$ and $\sqrt{5}$,
hence $\varphi$, and it is countable. On it the monotone cost forms are already
exactly the gauge orbit of $J$. So the cost form is forced on a countable carrier;
the cost layer does not require the continuum.

This does not erase the boundary, it relocates it. The number tower is forced. The
cost form is forced, now on a countable real-closed field rather than on the
uncountable completion. What still sits on the far side of a non-distinction posit
is only this: the genuinely transcendental physical constants the downstream theory
names, which use $\pi$ and the exponential and are therefore not algebraic. Even
those land in a countable field, obtained by adjoining their specific values to the
rationals, and that field is again a proper subset of the continuum. The continuum
is a scaffold for \emph{defining} the transcendental functions, not the home of any
answer the program returns. The honest residual posit is the existence of those
transcendental values, not the uncountable line itself.

\section{The meeting point}\label{sec:meeting}

We can now state the unification at exactly the strength the three layers support,
and not a word stronger.

\begin{theorem}[The stratified unification]\label{thm:unif}
Let $\delta$ be the act of distinction. Then:
\begin{enumerate}[leftmargin=2em, label=\textup{(\alph*)}]
  \item \textup{(Mathematics is what distinction permits.)} The number tower
    $\Ndelta \subset \Zdelta \subset \Qdelta$ and its permission structure are
    forced by $\delta$ under the distinction-only tag
    \textup{(Claims~\ref{cl:nat},~\ref{cl:rat})}.
  \item \textup{(The boundary.)} The continuous completion of $\Qdelta$ is not
    forced by $\delta$; it is a trace-closure commitment, separated from the
    distinction-only layer by a cardinality gap \textup{(Claim~\ref{cl:cont})}.
  \item \textup{(Physics is what the completion costs.)} On the completion, the
    cost requirements force the cost into the gauge orbit of $J$, with residual
    freedom exactly one positive real fixed by one calibration datum
    \textup{(Claims~\ref{cl:costneg},~\ref{cl:costpos})}.
\end{enumerate}
Mathematics and physics meet at one named boundary, the passage from the countable
ratio field to its uncountable completion, which is exactly the passage from the
distinction-only tag to the trace-closure tag.
\end{theorem}

\medskip
\noindent\textbf{The boundary on the cost side.} Part (c) is stated over the
completion to match the historical form of the uniqueness theorem, but
Section~\ref{sec:cost-monotone} shows the cost form is already forced over the
countable real-algebraic closure of $\Qdelta$, by order alone. The genuine
non-distinction commitment in part (b) is therefore needed only for naming the
transcendental constants, not for forcing the cost form. The meeting point is
sharper than the uncountable completion: it is the adjunction of specific
transcendental values to a countable field.

The meeting point is the heart of the program, and its character is worth stating
flatly. The cost is underdetermined on $\Qdelta$ for the same reason $\Qdelta$ is
incomplete: a gapped domain cannot transmit a local constraint, a curvature at the
unit, into a global one, a single cost function, because the gaps insulate the
constraint. The act of completing the field, removing the gaps, is the same act
that lets the four laws bite and force the cost form. So the boundary where the
mathematics ends and the physics begins is not two boundaries that happen to
align; it is one boundary seen from two sides. From the side of structure it is the
incompleteness of the ratio field. From the side of dynamics it is the
underdetermination of the cost. Completing the field closes both at once, at the
price of one posit.

This is why the program is stratified and not seamless, and why that is a strength.
A seamless chain from $\delta$ to the constants of physics would have to hide the
posit. The stratified account displays it, localizes it to one place, the
completion boundary, and one number, the calibration unit, and proves everything on
either side. The physics that the program develops downstream, the golden ratio and
the discrete time cadence and the spatial dimension and the numerical constants,
all live on the completed side of the boundary, and the flagship's task ends at the
boundary that makes them possible. We do not rehearse the downstream derivations
here; we locate them, correctly, as living beyond the posit, and we note that even
the first of them, the golden ratio, already requires the completion because
$\sqrt 5$ is not a ratio.

\section{Inevitability, stated conditionally}\label{sec:inev}

A natural further question is whether the distinction core is optional: could a
foundation be built that does not contain it? The program has a partial answer, and
because the answer is partial we state it as a conditional and resist the
temptation to inflate it.

\begin{claim}[Conditional inevitability]\label{cl:inev}
Once an admissible formal foundation is parsed into a common formal-system
interface, capturing its tokens, its formation rules, and its judgment of
syntactic identity, that foundation contains an embedding of the distinction core:
it has marks, it iterates them, and it judges them same or different, which is
exactly $\delta$ and its decidable comparison. In this sense the distinction core
is not an optional ingredient of such a foundation; it is presupposed by the act
of writing the foundation down.
\end{claim}

The honest qualification is essential. Claim~\ref{cl:inev} is proved for
foundations \emph{as parsed into the interface}. It is not yet proved that the major
historical foundations, set theory, type theory, category theory, and their
variants, each parse faithfully into the interface with their full expressive
power preserved. Until at least two such foundations are faithfully parsed and
their expressivity established, the inevitability claim is conditional on the
parsing, and to present it as unconditional would be to claim inevitability for our
own interface, which is circular. We therefore carry the inevitability material as
a conditional section of the synthesis rather than as a standalone theorem, and we
mark the parsing of real foundations as an open task. The defensible statement is
the modest one: any foundation, once written in a system of marks with a decidable
identity, already contains distinction, because writing it down is distinguishing.

\section{The honest accounting}\label{sec:accounting}

We collect the strength of every link in one place, so that the reader can see at a
glance what is forced, what is posited, and what is open. This table is the
program's ledger, and the discipline of keeping it is the program's principal
asset.

\begin{center}
\begin{tabular}{@{}lll@{}}
\toprule
\textbf{Link} & \textbf{Strength} & \textbf{Status} \\
\midrule
Natural numbers and their arithmetic & distinction-only & forced \\
Integers and their ring structure & distinction-only & forced \\
Rational field and its order & distinction-only & forced \\
Permission structure (divisibility, primality, gcd) & distinction-only & forced \\
Cost on the ratio field & distinction-only & not forced (classified) \\
Continuous completion of the ratio field & trace-closure & posited (not forced) \\
Cost form on the completion (gauge orbit of $J$) & trace-closure & forced \\
Selection of $J$ within the orbit & trace-closure & forced by one datum \\
Value of the calibration datum & trace-closure & posited (open: Section~\ref{sec:frontier}) \\
Inevitability of the distinction core & interface-conditional & conditional \\
Distinction as exactly one primitive & meta-level & open (Section~\ref{sec:frontier}) \\
\bottomrule
\end{tabular}
\end{center}

Read down the status column. The arithmetic layer is forced, entirely, under the
weakest tag. The continuum is posited, and provably so, not forced. The cost form
is forced, but only on the posited completion, and only up to a gauge. The gauge is
fixed by one datum, whose value is itself posited. Inevitability is conditional on
a parsing that is not yet done. And the question of whether distinction is one
primitive or two is open. Nothing in this ledger is dressed above its station, and
the two genuine open entries are named in the next section as research targets, not
hidden as solved.

\section{The open frontier}\label{sec:frontier}

Two questions are open, and one of them, if it closes, upgrades the program's
central claim. We state both precisely.

\subsection{The move that would force the cost outright}

The only free input standing between the imported continuous classification and an
unconditional forcing of $J$ is the calibration unit: the requirement that the
log-curvature of the cost at the unit equals one. On the completion this is
adopted, not derived. The target is to derive it from distinction directly.

\begin{claim}[Open: calibration from one distinction act]\label{cl:move1}
Conjecture: the minimal recognition cost of distinguishing one orbit from its
successor, the cost of a single act of distinction, carried to the completion,
fixes the second-order coefficient of the cost profile at the unit to exactly one.
If this is a theorem, the calibration datum is not a free choice but a consequence
of the primitive, the continuous classification runs with no free input, and
``distinction forces $J$ on the completion'' becomes true without circularity. If
it is false, the genuine native curvature is some other value, the forced cost is
the corresponding member of the gauge orbit rather than $J$, and one has learned
the true cost of distinction, which is itself a result.
\end{claim}

We do not claim either outcome. We mark the question as the single
highest-leverage open problem of the program, with a falsifier built in: the
one-act cost expansion either has leading curvature one or it does not, and both
answers are publishable and more fundamental than the present conditional. The
program is built to ship its current, honest claim without this result and to
absorb it cleanly if it arrives.

\subsection{One primitive or two}

The second open question is the one raised in Section~\ref{sec:primitive}: whether
the same-or-different judgment is derivable from the act of marking, in which case
distinction is genuinely one primitive, or whether recognition is irreducibly the
pair, mark and compare. Resolving it does not change any downstream result, since
both moves are available throughout, but it changes what the program may honestly
boast. Either the judgment reduces to the act, and the ``single primitive'' claim
is earned, or it does not, and the honest description is that recognition is the
pair. Both are more fundamental than asserting unity without proof, and we leave
the question open rather than assert the slogan.

\section{Conclusion}

The program takes one primitive, the act of distinction, and the honest measure of
its reach is a stratification with one boundary. Distinction forces the number
tower, all of it, under the weakest logical commitments, with no infinity, no
choice, and no excluded middle: mathematics is what distinction permits.
Distinction does not force the continuous line that physics requires; the line is a
strictly larger object, separated from distinction's countable reach by a gap in
cardinality, and adopted as a named posit. On that posited completion, and only
there, the cost is forced into the gauge orbit of the canonical reciprocal cost,
with exactly one real degree of freedom fixed by one calibration datum: physics is
what the completion of distinction costs. The two layers meet at one boundary, the
completion of the ratio field, which is simultaneously the end of the
distinction-only tag and the beginning of the trace-closure tag, the end of
mathematics-as-permitted and the beginning of physics-as-costed.

The strength of this account is precisely its refusal to overstate. It localizes
the program's entire posited content to one boundary and one number, proves
everything on both sides, treats the inevitability of the core as the conditional
it currently is, and names the two open questions as research targets rather than
solved problems. A grand claim that arrives with its boundary in the abstract and
its open frontier named in the text is a different and sturdier object than a
seamless chain that hides its joints. The unification is real, and it is
stratified. Distinction permits mathematics; the completion of distinction costs
physics; and the single open move that would close the last gap, deriving the unit
of cost from the cost of one act of distinction, is stated here as the frontier it
is.

\begin{thebibliography}{99}

\bibitem{DeltaArithmetic}
J.~Washburn,
``Arithmetic from Distinction: A Number Tower Forced by a Single Primitive,''
companion paper, 2026.

\bibitem{DeltaContinuum}
J.~Washburn,
``Where Necessity Ends: Why Distinction Forces Arithmetic but Not the
Continuum,'' companion paper, 2026.

\bibitem{DeltaCost}
J.~Washburn,
``The Cost Form Is Forced, the Unit Is a Gauge: How Far Distinction Determines
the Recognition Cost,'' companion paper, 2026.

\bibitem{LogicFE}
J.~Washburn,
``A Functional-Equation Encoding of Logical Consistency on Continuous
Positive-Ratio Comparisons: A Conditional Rigidity Theorem,'' 2026.

\bibitem{AxiomsCost}
J.~Washburn and Z.~Zlatanovi\'c,
``Uniqueness of the Canonical Reciprocal Cost,''
\emph{Axioms}, 2026.

\bibitem{Aczel1966}
J.~Acz\'el,
\emph{Lectures on Functional Equations and Their Applications},
Academic Press, 1966.

\bibitem{Cantor1891}
G.~Cantor,
``\"Uber eine elementare Frage der Mannigfaltigkeitslehre,''
\emph{Jahresbericht der Deutschen Mathematiker-Vereinigung} \textbf{1} (1891),
75--78.

\end{thebibliography}

\end{document}
